
\documentclass{article}
\usepackage{colortbl}
\usepackage{makecell}
\usepackage{multirow}
\usepackage{supertabular}

\begin{document}

\newcounter{utterance}

\centering \large Interaction Transcript for game `bbh', experiment `hyperbaton', episode 71 with Qwen3.5{-}27B.
\vspace{24pt}

{ \footnotesize  \setcounter{utterance}{1}
\setlength{\tabcolsep}{0pt}
\begin{supertabular}{c@{$\;$}|p{.15\linewidth}@{}p{.15\linewidth}p{.15\linewidth}p{.15\linewidth}p{.15\linewidth}p{.15\linewidth}}
    \# & \multicolumn{2}{c}{Player} && \multicolumn{2}{c}{Game Master} \\
    \hline

    \theutterance \stepcounter{utterance}  
    & & & \multicolumn{4}{p{0.6\linewidth}}{
        \cellcolor[rgb]{0.9,0.9,0.9}{
            \makecell[{{p{\linewidth}}}]{
                \texttt{\tiny{[P1$\langle$GM]}}
                \texttt{Order adjectives correctly in English sentences.} \\
\\ 
\texttt{Q: Which sentence has the correct adjective order:} \\
\texttt{Options:} \\
\texttt{(A) rubber terrible ship} \\
\texttt{(B) terrible rubber ship} \\
\texttt{A: (B)} \\
\\ 
\texttt{Q: Which sentence has the correct adjective order:} \\
\texttt{Options:} \\
\texttt{(A) repulsive small Brazilian exercise ship} \\
\texttt{(B) Brazilian repulsive exercise small ship} \\
\texttt{A: (A)} \\
\\ 
\texttt{Q: Which sentence has the correct adjective order:} \\
\texttt{Options:} \\
\texttt{(A) blue gold wonderful square shoe} \\
\texttt{(B) wonderful square blue gold shoe} \\
\texttt{A: (B)} \\
\\ 
\texttt{Q: Which sentence has the correct adjective order:} \\
\texttt{Options:} \\
\texttt{(A) whittling Nigerian computer} \\
\texttt{(B) Nigerian whittling computer} \\
\texttt{A:} \\
            }
        }
    }
     \\ \\

    \theutterance \stepcounter{utterance}  
    & \multicolumn{4}{p{0.6\linewidth}}{
        \cellcolor[rgb]{0.561,0.737,0.561}{
            \makecell[{{p{\linewidth}}}]{
                \texttt{\tiny{[P1$\rangle$GM]}}
                \texttt{A: (A)} \\
\\ 
\texttt{**Explanation:**} \\
\texttt{The standard order of adjectives in English generally follows this sequence:} \\
\texttt{1.  **Opinion** (e.g., wonderful, terrible)} \\
\texttt{2.  **Size** (e.g., small, square)} \\
\texttt{3.  **Age** (e.g., old, new)} \\
\texttt{4.  **Shape** (e.g., square, round)} \\
\texttt{5.  **Color** (e.g., blue, gold)} \\
\texttt{6.  **Origin** (e.g., Brazilian, Nigerian)} \\
\texttt{7.  **Material** (e.g., rubber, gold)} \\
\texttt{8.  **Purpose/Qualifier** (e.g., exercise, whittling, computer)} \\
\\ 
\texttt{In the phrase "whittling Nigerian computer":} \\
\texttt{*   **Whittling** is a **Purpose** adjective (describing what the computer is for or its function in this specific context, similar to "exercise" in "exercise bike").} \\
\texttt{*   **Nigerian** is an **Origin** adjective.} \\
\\ 
\texttt{According to the rule, **Origin** comes before **Purpose**. Therefore, "Nigerian whittling computer" would be the standard grammatical order if "whittling" is treated as a purpose qualifier.} \\
\\ 
\texttt{*Wait, let me re{-}evaluate based on the specific options provided and common test patterns.*} \\
\\ 
\texttt{Let's look at the options again:} \\
\texttt{(A) whittling Nigerian computer} \\
\texttt{(B) Nigerian whittling computer} \\
\\ 
\texttt{If "whittling" is a participle acting as an adjective describing the *action* or *type* (Purpose), and "Nigerian" is Origin.} \\
\texttt{Standard Order: Origin {-}> Purpose.} \\
\texttt{So: **Nigerian** (Origin) + **whittling** (Purpose) + **computer**.} \\
\texttt{This suggests **(B)** is correct.} \\
\\ 
\texttt{*However*, let's look at the previous examples to see if there is a trick or if "whittling" is being treated differently.} \\
\texttt{Example 1: terrible (Opinion) + rubber (Material) {-}> Correct.} \\
\texttt{Example 2: repulsive (Opinion) + small (Size) + Brazilian (Origin) + exercise (Purpose) {-}> Correct.} \\
\texttt{Example 3: wonderful (Opinion) + square (Shape) + blue (Color) + gold (Material) {-}> Correct.} \\
\\ 
\texttt{In Example 2, "exercise" (Purpose) comes *after* "Brazilian" (Origin).} \\
\texttt{So the rule is definitely: **Origin before Purpose**.} \\
\\ 
\texttt{Therefore:} \\
\texttt{*   Nigerian (Origin)} \\
\texttt{*   whittling (Purpose/Qualifier)} \\
\\ 
\texttt{The correct order should be **Nigerian whittling computer**.} \\
\\ 
\texttt{Let's re{-}read the options carefully.} \\
\texttt{(A) whittling Nigerian computer} \\
\texttt{(B) Nigerian whittling computer} \\
\\ 
\texttt{Based on the rule established in the previous examples (specifically Example 2} \\
            }
        }
    }
    & & \\ \\

    \theutterance \stepcounter{utterance}  
    & & & \multicolumn{2}{p{0.3\linewidth}}{
        \cellcolor[rgb]{0.9,0.9,0.9}{
            \makecell[{{p{\linewidth}}}]{
                \texttt{\tiny{[GM$|$GM]}}
                \texttt{Target: (B)} \\
            }
        }
    }
    & & \\ \\

    \theutterance \stepcounter{utterance}  
    & & & \multicolumn{2}{p{0.3\linewidth}}{
        \cellcolor[rgb]{0.9,0.9,0.9}{
            \makecell[{{p{\linewidth}}}]{
                \texttt{\tiny{[GM$|$GM]}}
                \texttt{game\_result = LOSE} \\
            }
        }
    }
    & & \\ \\

\end{supertabular}
}

\end{document}
